\section{Real realizations and moment-graph sheaves}
\label{sec:background}

We first construct a real realization to which the character theorems apply.
We then state the reflection-order and sheaf inputs, ending with the two graded ranks
that will be the endpoints of the comparison in Section~\ref{sec:induction}.
All constructions are made separately for the two Coxeter systems.
All sheaf constructions use the coefficient field $\mathbb R$, and
linear forms have degree one.
Theorem and section locators for \cite{Fiebig08,EW14} refer to the
arXiv versions specified in the bibliography.

\subsection{Reduction to finite rank}

Let $J$ be the finite set of simple reflections occurring in a reduced
expression for $b$. Every element below $b$ lies in the standard parabolic
subgroup $W_J$, by the subword criterion for Bruhat order
\cite[Theorem~2.2.2]{BB05}. The Bruhat order on $W_J$ is the restriction of
the ambient order \cite[Proposition~2.4.1]{BB05}, and its equal-parameter
Hecke algebra is the corresponding standard subalgebra. The defining
recursions therefore give the same $R$- and Kazhdan--Lusztig polynomials on
this subgroup. The interval graph also restricts. For an ambient edge
$x\to y$ with $x,y\in W_J$, its label $t=yx^{-1}$ belongs to $W_J$.
Strong exchange applied to $tt=1$ and a reduced $J$-word for $t$ expresses
$t$ as a $W_J$-conjugate of a generator in $J$
\cite[Corollary~1.4.4]{BB05}; the converse is immediate.
Consequently we may, and do,
assume that both Coxeter systems have finite rank.  The groups themselves
need not be finite.

\subsection{A reflection-faithful real realization}

\begin{lemma}
\label{bg:realization}
For a finite-rank Coxeter system $(W,S)$ there is a finite-dimensional real
vector space $V$, with a nondegenerate symmetric bilinear form, having the
following properties.
\begin{enumerate}
\item The simple roots $\alpha_s$ are linearly independent,
\[
 (\alpha_s,\alpha_s)=2,
 \qquad
 (\alpha_s,\alpha_t)=-2\cos(\pi/m_{st})\quad(s\ne t),
\]
where $\cos(\pi/\infty)=1$, and $s$ acts by
\[
 s(v)=v-(\alpha_s,v)\alpha_s.
\]
\item The representation is reflection faithful: it is faithful, and
$\operatorname{rank}(w-1)=1$ holds precisely for the reflections of $W$.
\item After identifying $V$ with its dual by the bilinear form, both the
simple roots and the simple coroots are independent.  A strictly dominant
linear form $\rho$ can be chosen so that
\[
 (w\rho)(\alpha_s)>0\quad\Longleftrightarrow\quad sw>w.
\]
\end{enumerate}
This realization may be chosen to have the minimal dimension required by
the independent-root and independent-coroot conditions in
\cite[Section~3.1]{EW14}.
\end{lemma}

\begin{proof}
Start with the usual geometric root space $V_0$ and its symmetric Tits form
\cite[Sections~4.2--4.4]{BB05}.
Let $K$ be its radical and choose a complement $V_1$ to $K$.  Adjoin a vector
space $K'$ dual to $K$, pair $K$ and $K'$ nondegenerately, and take $K'$ to be
isotropic and orthogonal to $V_1$.  The resulting form on
\[
 V=V_1\oplus K\oplus K'
\]
is nondegenerate and extends the original form.  Its dimension is
$\dim V_0+\dim K$, the minimal possible dimension of such an extension.
This is also minimal under the independent-root and independent-coroot
conditions.  Indeed, let $G=((\alpha_s,\alpha_t))_{s,t\in S}$.  In any
realization $\widetilde V$ with these pairings and independent simple roots
and coroots, evaluation by the root covectors maps $\widetilde V$ onto $\R^S$.
Its restriction to the coroot span has matrix $G$, so its kernel there
has dimension $\operatorname{corank}G$.  This subspace lies in the full
kernel, so rank--nullity gives
\[
 \dim\widetilde V\ge |S|+\operatorname{corank}G=\dim V_0+\dim K,
\]
which the construction attains.
For distinct simple reflections with $m_{st}<\infty$, their root plane is
positive definite. The displayed reflections have product of order $m_{st}$
on that plane and fix its orthogonal complement. Thus the Coxeter relations
hold on all of $V$; pairs with $m_{st}=\infty$ impose no relation.
The action restricts to the faithful geometric action on $V_0$
\cite[Theorem~4.2.7]{BB05},
so it is faithful on $V$ as well.
The simple roots and their corresponding linear forms are independent by
nondegeneracy.

For clarity, reflection faithfulness can also be seen directly.  If an
orthogonal transformation has rank-one difference from the identity, write
it as $g=1+u\otimes\varphi$.  Orthogonality first implies
$u^{\flat}=c\varphi$ for some nonzero scalar $c$, and then implies
$\varphi(u)=-2$.  Thus $g^2=1$ and $(u,u)\ne0$: $g$ is an orthogonal
reflection.  If $g\in W$, the finite-subgroup theorem for Coxeter groups
\cite[Theorem~12.3.4(i)]{Davis08} puts
$\langle g\rangle$, after conjugation, in a finite standard parabolic
subgroup.  Its root span is positive definite, and its orthogonal complement
in $V$ is fixed by that parabolic subgroup.  The restriction of $g$ to the
root span still has rank-one difference from the identity, and finite
Coxeter geometry identifies it as a Coxeter reflection.  Conversely every
Coxeter reflection has rank-one difference from the identity.  This is the
reflection-faithful realization used in \cite[Section~3.1]{EW14}.

Choose $\rho\in V^*$ positive on every simple coroot.  Such a choice exists
by their independence.  A real coroot has either nonnegative or nonpositive
simple-coroot coefficients.  Applying this observation to
$w^{-1}\alpha_s$ gives the displayed length-sign criterion.  The same
argument applies to roots under the symmetric identification.

The chosen $\rho$ has trivial stabilizer. If $w\ne1$, take a left descent
$s$ of $w$; then $(w\rho)(\alpha_s)<0<\rho(\alpha_s)$, so $w\rho\ne\rho$.
The identification of $V$ with its dual therefore supplies a regular orbit.
The comparison constructions of \cite[Sections~2 and~5]{Fiebig08} also
use an orbit with stabilizer $\{1,s\}$ for each simple reflection $s$.
Such an orbit exists here. Put $H_s=V^s$. An element fixing $H_s$ pointwise
has rank at most one after subtracting the identity, so reflection
faithfulness and uniqueness of the orthogonal reflection with hyperplane
$H_s$ force it to be $1$ or $s$. For every other element $w$,
$H_s\cap V^w$ is therefore a proper subspace of $H_s$.
Since $W$ is countable, a point outside their union exists over $\R$
and has stabilizer exactly $\{1,s\}$.
\end{proof}

Let $\Phi=\{w\alpha_s:w\in W,\ s\in S\}$ be the real root system,
and let $\Phi^+$ be its positive roots. For $\gamma\in\Phi^+$, write
$r_\gamma$ for its reflection. The extension leaves the root action
unchanged, so the standard reflection length-sign criterion
\cite[Proposition~4.4.6]{BB05}, applied to $w^{-1}$, gives
\begin{equation}\label{eq:bg-reflection-sign}
 r_\gamma w>w
 \quad\Longleftrightarrow\quad
 \ell(r_\gamma w)>\ell(w)
 \quad\Longleftrightarrow\quad
 w^{-1}\gamma\in\Phi^+.
\end{equation}
Put
\[
 A=\operatorname{Sym}_{\mathbb R}(V),\qquad N=\dim V,
\]
where the bilinear form identifies each root with a linear polynomial.
Throughout the proof, linear polynomials have degree $1$.  If $E$ is an edge
labeled by a reflection $t$, write $\alpha_E$ for any nonzero scalar multiple
of its root.  Only the principal ideal $(\alpha_E)$ matters.  Distinct
reflections give nonproportional labels.  We include the additional
directions of $V$ in the polynomial ring, even when the original Tits form
is degenerate.

For a finitely generated graded free $A$-module $M$ with homogeneous basis
degrees $j_1,\ldots,j_r$, set
\[
 \operatorname{grk}_A(M;q)=\sum_{i=1}^r q^{j_i}.
\]
Thus a generator in positive degree contributes a positive exponent, and
\[
 \operatorname{Hilb}(M;q)
 =\frac{\operatorname{grk}_A(M;q)}{(1-q)^N}.
\]

\subsection{Reflection orders and their path formula}

A reflection order is a total order of the positive roots whose
restriction to the positive roots in every rank-two root subsystem is one
of their two angular orders.  We shall use geometric reflection orders.
Choose a linear functional $h$ strictly positive on the simple roots, and
therefore on every positive root.  Choose another linear functional $f$
such that the numbers
\[
 \frac{f(\alpha)}{h(\alpha)},\qquad \alpha\in\Phi^+,
\]
are pairwise distinct, and order the roots by these numbers.  Such an $f$
exists: the root set is countable, and each equality to be avoided cuts
out a proper hyperplane in the space of functionals.  Within a root plane,
the normalized positive roots $\alpha/h(\alpha)$ lie on an affine line;
their order is consequently one of the angular orders.  This produces a
reflection order even when the root system is infinite.  Reversing it
produces another reflection order.

\begin{theorem}[Dyer's reflection-order path formula]
\label{bg:paths}
Fix a reflection order and use left reflection labels on Bruhat edges.
Let $\widetilde R_{x,y}(T)$ be the sum of $T^r$ over directed paths of $r$
edges from $x$ to $y$ having strictly increasing labels.  The empty path
contributes $1$ when $x=y$.  Then
\begin{equation}
\label{eq:bg-path-formula}
 q^{-d(x,y)/2}R_{x,y}(q)
 =\widetilde R_{x,y}\bigl(q^{1/2}-q^{-1/2}\bigr).
\end{equation}
The same polynomial counts paths with strictly decreasing labels.
\end{theorem}

The path formula is due to Dyer \cite[Corollary~3.4]{Dyer93}; its left-label version
follows from the right-label version by inversion.  The decreasing
statement follows by reversing the reflection order.  Every path between
$x$ and $y$ lies in the finite interval $[x,y]$, so all these sums are
finite.  Finally, edges with the same reflection label are pairwise
disjoint: a reflection pairs each vertex $w$ only with $tw$.  Thus any
schedule obtained by ordering reflection labels has no interacting ties.

\subsection{The sheaf and its full upward kernels}

Fix $b\in W$.  The Bruhat graph on $[1,b]$ has an edge
$E:x\longrightarrow y$ whenever $y=tx>x$ for a reflection $t$; labels are
on the left.  A sheaf on this graph consists of vertex modules $B^x$, edge
modules $B^E$, and restriction maps from the two endpoints of each edge.
A section over a set of vertices is a tuple whose endpoint restrictions
agree on every edge with both endpoints in the set.

We use the following standard form of the Braden--MacPherson construction
and its relation to Soergel bimodules.  The construction and the Verma-flag
properties are due to Braden--MacPherson and Fiebig
\cite{BM01,FiebigVerma08,Fiebig08}, within Soergel's categorification
\cite{Soergel07}; the character identity
used below is the theorem of Elias--Williamson.  The precise bridges are
\cite[Theorem~6.3, Corollary~6.5, and Proposition~7.1]{Fiebig08},
\cite[Section~3]{FiebigSurvey}, and \cite[Theorem~1.1]{EW14}.

\begin{theorem}[Moment-graph sheaf]
\label{bg:bmp}
There is an indecomposable graded sheaf $B=B(b)$ on $[1,b]$, normalized by
$B^b=A$, with the following properties.
\begin{enumerate}
\item Every vertex module $B^x$ is finitely generated and graded free.
For $E:x\longrightarrow y$,
\[
 B^E=B^y/\alpha_E B^y,
\]
and restriction from $y$ is the quotient map.
\item If $\Omega\subseteq[1,b]$ is upward closed, every section over
$\Omega$ extends to a global section.  Moreover the projection from global
sections onto each $B^x$ is surjective.
\item Let $U_x$ be the set of edges out of $x$.  For $F\subseteq U_x$ put
\[
 L_x(F)=\ker\left(B^x\longrightarrow
                   \bigoplus_{E\in F}B^E\right).
\]
Then $L_x(U_x)$ is graded free.
\end{enumerate}
\end{theorem}

The second property will be called flabbiness and generation by global
sections.  No freeness assertion for an arbitrary subset $F\subseteq U_x$
is included in this theorem.  Freeness for the particular subsets needed
later will be proved by induction.

\begin{proposition}[Grading normalization]
\label{bg:normalization}
For $x\le b$, write $d=d(x,b)=\ell(b)-\ell(x)$.  In the degree-one
convention above,
\begin{equation}
\label{eq:bg-ranks}
 \operatorname{grk}_A(B^x;q)=P_{x,b}(q),
 \qquad
 \operatorname{grk}_A(L_x(U_x);q)=q^dP_{x,b}(q^{-1}).
\end{equation}
\end{proposition}

\begin{proof}
We track the shifts explicitly.  Temporarily double all degrees, so that
linear forms have degree $2$, and let $z$ record the resulting generator
degrees.  Denote the regraded sheaf by $\widehat B$, and write
$n=\ell(b)$ and $m=\ell(x)$.  With the shift convention
$M(k)^j=M^{j+k}$, put
\[
 M=\Gamma(\widehat B)(n)
\]
and note that its top generator has degree $-n$.
This is the normalized self-dual indecomposable projective in
\cite[Corollary~6.5 and Theorem~6.1]{Fiebig08}.

We first compute the full upward-kernel rank.  For $w\le b$, let
$M^{[w]}$ be the lower-support layer
\[
 M^{[w]}=\Gamma_{\le w}M/\Gamma_{<w}M,
\]
where $\Gamma_JM$ means the submodule of sections supported on $J$.
Restriction to $x$, together with flabbiness, identifies $M^{[x]}$ with
$\widehat L_x(U_x)(n)$.  Indeed a vector in the full upward kernel, placed
at $x$ and extended by zero outside $\{y:y\le x\}$, is a section on the
upward-closed set $[1,b]\setminus\{y:y<x\}$; flabbiness extends it to a
global section supported on $\{y:y\le x\}$.

The layer in \cite[Section~2.4]{Fiebig08} is described instead as the
kernel of the map between the projections to the upper sets
$\{y:y\ge x\}$ and $\{y:y>x\}$. Flabbiness identifies those projections
with section modules; a tuple in the kernel has only its $x$-coordinate
nonzero, with precisely the full upward vanishing conditions.
Thus that layer and the lower-support quotient above are canonically
identified through $\widehat L_x(U_x)(n)$.

We now compare the character conventions in the cited sources.  In the
Hecke normalization $H_s^2=1+(z^{-1}-z)H_s$, write $\underline H_w$ for
the Kazhdan--Lusztig basis element.  The Hecke bar involution is determined
by $\overline z=z^{-1}$ and $\overline{H_s}=H_s^{-1}$.
Fiebig's twisted graph indexed by $w$
is $\{(w^{-1}v,v):v\in V\}$ \cite[Section~4]{Fiebig08},
whereas Elias--Williamson use $\{(wv,v):v\in V\}$
\cite[Section~3.5]{EW14}.
Under Fiebig's equivalence with Soergel bimodules
\cite[Theorem~6.3]{Fiebig08}, this index reversal and
the top normalization identify the corresponding bimodule as the unshifted
$B_{b^{-1}}$.  Let $\mathfrak a$ be the $\Z[z,z^{-1}]$-linear
anti-involution of the Hecke algebra defined by
$\mathfrak a(H_w)=H_{w^{-1}}$ \cite[Section~3.2]{EW14}.

Fiebig's lower-support character of $M$ is
$\mathfrak a(\operatorname{ch}_{\nabla}(B_{b^{-1}}))$ in
Elias--Williamson's notation.  The coefficient convention can be checked
on one generator.  Their standard lower-support module
$\nabla_{x^{-1}}$ has its generator in degree $-m$.  A generator of
$M^{[x]}$ in degree $j$ therefore corresponds to
$\nabla_{x^{-1}}(-m-j)$.  Its shift contributes $z^{-m-j}$ before the
coefficient conjugation in $\operatorname{ch}_{\nabla}$ changes this
to $z^{m+j}$.  Thus the coefficient after applying $\mathfrak a$ is
$z^m\operatorname{grk}_z(M^{[x]})$, as in
\cite[Section~8.2]{Fiebig08}.

Let $\operatorname{ch}$ denote the standard character in
\cite[Theorem~1.1]{EW14}.  For a Soergel bimodule, the characters satisfy
$\operatorname{ch}_{\nabla}(B)=\overline{\operatorname{ch}(B)}$
\cite[Section~3.5]{EW14}.  The character theorem, together with bar invariance of the
Kazhdan--Lusztig basis, therefore gives
$\operatorname{ch}_{\nabla}(B_{b^{-1}})=\underline H_{b^{-1}}$.
The anti-involution $\mathfrak a$ commutes with bar and inversion preserves
Bruhat order, so the triangular characterization gives
$\mathfrak a(\underline H_{b^{-1}})=\underline H_b$.
Using the Kazhdan--Lusztig expansion in \cite[Remark~3.2]{EW14}, the full
character identity is
\begin{align*}
 \sum_{w\le b}z^{\ell(w)}\operatorname{grk}_z(M^{[w]})H_w
 &=\mathfrak a(\operatorname{ch}_{\nabla}(B_{b^{-1}}))
   =\underline H_b\\
 &=\sum_{w\le b}z^{n-\ell(w)}P_{w,b}(z^{-2})H_w.
\end{align*}
Comparison of the coefficient of $H_x$ gives
\begin{equation}
\label{eq:bg-character}
 z^m\operatorname{grk}_z(M^{[x]})
   =z^{n-m}P_{x,b}(z^{-2}).
\end{equation}
Since $M^{[x]}\simeq\widehat L_x(U_x)(n)$, this gives
\begin{equation}
\label{eq:bg-upward-rank}
 \operatorname{grk}_z\widehat L_x(U_x)
 =z^{2(n-m)}P_{x,b}(z^{-2}).
\end{equation}

For the stalk rank, write $M^x$ for the image of the projection of $M$
to $\widehat B^x(n)$.  Generation by global sections gives
$M^x=\widehat B^x(n)$.
Fiebig's integral graded-module statement
\cite[Proposition~7.1(4), with Corollary~6.5]{Fiebig08} says, in the
shift convention above, that
\[
 M^x\simeq\bigoplus_i A(k_i)
 \quad\Longrightarrow\quad
 M^{[x]}\simeq\bigoplus_i A(2m-k_i).
\]
Both sides are the actual graded free modules over $A$.  In our
positive-generator rank convention, this relation becomes
\[
 \operatorname{grk}_z(M^x)
 =z^{-2m}\operatorname{grk}_z(M^{[x]})\big|_{z\mapsto z^{-1}}.
\]
Using $M^x=\widehat B^x(n)$ and \eqref{eq:bg-character}, we obtain
\[
 z^{-n}\operatorname{grk}_z(\widehat B^x)
 =z^{-2m}\left.z^{n-2m}P_{x,b}(z^{-2})\right|_{z\mapsto z^{-1}}
 =z^{-n}P_{x,b}(z^2).
\]
Thus $\operatorname{grk}_z\widehat B^x=P_{x,b}(z^2)$.
Finally set $q=z^2$ in this formula and in
\eqref{eq:bg-upward-rank}.  This gives \eqref{eq:bg-ranks}.
\end{proof}

\begin{remark}
\label{bg:upper-interval}
Although $B(b)$ is constructed on $[1,b]$, every upward edge out of a
vertex $x\in[u,b]$ has its other endpoint in $[u,b]$.  Thus $U_x$ and all
the modules $L_x(F)$ used on this upper interval are computed without any
additional upward conditions outside it.
\end{remark}
